% =========================================================
\section{Implementation Details and Model Prompts}
\label{app:implementation}
% =========================================================

\subsection{Rule-Based Set Score for First-Round Hard Negatives}
\label{app:hand_score}

Because no learned set ranker is available before the first RankNet training round, we first use a fixed rule-based set score
$H(\Selected)$ to select harder negative samples from incomplete candidate sets. This score is used only for negative-sample selection and is not the final TEF-RAG prediction score. It is defined as
\begin{equation}
\begin{aligned}
H(\Selected)
={}&S_{\mathrm{node}}(\Selected)
+0.48\,S_{\mathrm{edge}}(\Selected)
+0.42\,S_{\mathrm{role}}(\Selected)\\
&-0.16\,S_{\mathrm{red}}(\Selected)
-0.10\,S_{\mathrm{disc}}(\Selected).
\end{aligned}
\label{eq:appendix_hand_score}
\end{equation}

The node term is the sum of the evidence-node scores within the set:
\begin{equation}
S_{\mathrm{node}}(\Selected)
=
\sum_{e\in\Selected}s_{\mathrm{node}}(e),
\end{equation}
where the score of an individual evidence record is
\begin{align}
s_{\mathrm{node}}(e)
=
0.62\,s_{\mathrm{rel}}(e)
+0.18\,s_{\mathrm{rec}}(e)
+0.15\,s_{\mathrm{role}}(e) \\
+0.05\,s_{\mathrm{src}}(e). \notag
\end{align}
Here, $s_{\mathrm{rel}}$ is the min--max-normalized BM25 score among the current legal candidates.
The temporal-recency term is $s_{\mathrm{rec}}=\exp(-\Delta t/T)$, with $T=45$ days for ordinary queries and $T=180$ days for queries involving comparisons across time.
If the evidence type matches a task role required by the current query, $s_{\mathrm{role}}=1$; otherwise it is set to $0.35$.
In the current implementation, $s_{\mathrm{src}}=1$ for all legal candidates, so this term contributes only a constant offset and does not affect ranking among evidence nodes.

The relation term counts only valid relation edges internal to the candidate set. To avoid giving an unfair advantage to dense graphs simply because they contain more edges, only the highest-confidence incoming edge is retained for each target node:
\begin{equation}
S_{\mathrm{edge}}(\Selected)
=
\sum_{e_j\in\Selected}
\max_{(e_i,r,e_j)\in E(\Selected)}c_{ij},
\end{equation}
where a node with no incoming edge contributes 0.

Role coverage is defined as
\begin{equation}
S_{\mathrm{role}}(\Selected)
=
\frac{|R(\Selected)\cap D_q|}{|D_q|},
\end{equation}
where $R(\Selected)$ is the set of evidence roles appearing in the candidate set and $D_q$ is the set of roles required by the current query. The system considers state observation, diagnosis, and work order by default, and adds roles such as procedure, verification, inspection, uncertainty, or correction when indicated by the query.

The redundancy term is based on Jaccard similarity between token sets of evidence texts:
\begin{equation}
S_{\mathrm{red}}(\Selected)
=
\sum_{i<j}
\max\left(
0,
\operatorname{Jaccard}(e_i,e_j)-0.45
\right).
\end{equation}
The isolation term $S_{\mathrm{disc}}$ counts evidence records in the set that do not participate in any valid relation edge.
During the first training round, candidate sets with high $H(\Selected)$ but incomplete reference evidence flows are prioritized as negative samples. The second round no longer relies on this rule-based score; instead, incomplete sets receiving the highest scores from the current RankNet are selected directly as hard negatives.

\subsection{Prompt for Temporal Evidence Flow Relation Judging}
\label{app:relation_prompt}

The relation judge uses a fixed system prompt (version \texttt{tef-v6-stage2a-relation-v7}).
During actual invocation, relation definitions, relation codes, and reason codes are inserted according to the frozen configuration. The system prompt is:

\begin{lstlisting}
Prompt version: tef-v6-stage2a-relation-v7. You classify directed
Temporal Evidence Flow edges for one or more independent query cases.
Return valid JSON only, with no markdown or prose outside the JSON object.
Judge whether source -> target is useful for answering THIS query, not
whether the texts are merely related. Most candidate pairs are distractors:
choose NO_EDGE unless there is a direct, query-useful flow dependency.
Shared asset, episode, topic, or chronology alone never licenses an edge.
Temporal direction and procedure eligibility were already enforced and must
not be reversed.

Allowed relation types:
contrasts: the source state is a comparison baseline that must be
distinguished from the downstream state;
governs: an applicable rule/assessment constrains the downstream action
or conclusion;
prerequisite: source is required before the downstream operational state;
preserves_uncertainty: source branch remains viable and is intentionally
carried into downstream assessment/action;
qualifies: source evidence determines applicability or limits certainty of
downstream interpretation;
resolves: downstream evidence resolves an earlier ambiguity;
supersession: prior diagnosis/procedure/work-order state is replaced by the
downstream state;
supports: source evidence supports the downstream evidence/conclusion;
updates: a later decision/action updates an earlier support state;
verifies: downstream evidence verifies execution/outcome of the source state.
Otherwise use NO_EDGE.

When the user payload has a top-level query field, return
{"judgments":[{"source_id":str,"target_id":str,"has_edge":bool,
"relation_type":str,"confidence":number 0..1,
"reason_code":short_snake_case}]}.

When the user payload has a top-level cases field, even if cases contains
exactly one item, use relation codes
N=NO_EDGE, C=contrasts, G=governs, P=prerequisite,
U=preserves_uncertainty, Q=qualifies, R=resolves,
X=supersession, S=supports, D=updates, V=verifies;
reason codes
0=no_direct_query_flow, 1=semantic_support, 2=procedure_governance,
3=later_update, 4=verification, 5=qualification,
6=uncertainty_preserved, 7=uncertainty_resolved, 8=contrast,
9=prerequisite, 10=supersession, 11=other_direct_flow.

Pairs are ordered. Return
{"cases":[[case_id,[[relation_code,confidence_percent,reason_code],...]],...]}.
Each inner result list must have exactly the same length and order as that
case's supplied pairs. The user supplies deduplicated evidence tables and
source_id/target_id pairs. Do not mix facts across cases. Do not repeat
evidence IDs. Do not provide prose or hidden reasoning.
\end{lstlisting}

The corresponding user input contains only the query, candidate evidence, and evidence pairs to be judged, with the following structure:

\begin{lstlisting}
{
  "query": {
    "query_text": ...,
    "query_time": ...,
    "asset_model": ...,
    "asset_context": ...
  },
  "evidence": [...],
  "pairs": [
    {"source_id": ..., "target_id": ...},
    ...
  ]
}
\end{lstlisting}

\subsection{Prompt for Downstream Structured Generation}
\label{app:generation_prompt}

All retrieval methods share the same frozen downstream generator. Its system prompt is:

\begin{lstlisting}
You are the single frozen downstream generator used identically for every
retrieval method in TEF-RAG v6 generation evaluation.

Produce ONLY a JSON object matching the supplied schema. You receive one
query and exactly the evidence selected by a retrieval method. The retrieval
method identity is hidden and irrelevant.

Rules:
- Use only the supplied evidence. Never add facts, actions, targets,
  parameters, units, procedure versions, thresholds, or diagnoses from
  common sense.
- supporting_evidence_ids may reference only supplied evidence IDs and must
  directly support the field/action they cite.
- work_order.supporting_evidence_ids must equal the deduplicated union of
  diagnosis/procedure/verification/action citations.
- asset_id must equal the query asset_id.
- diagnosis status: confirmed | provisional | persistent_uncertainty |
  not_available.
- procedure status: applicable | superseded | uncertain | not_applicable.
  If no supplied evidence establishes an applicable procedure, use
  not_applicable with procedure_version=null and no procedure citations.
- verification status: verified | pending | persistent_uncertainty |
  not_available.
- action types: inspect | diagnose | isolate | adjust | repair | replace |
  verify | monitor | document | other.
- action_plan order in the JSON array has no semantic meaning. Use depends_on
  only for necessary prerequisite edges; keep the graph acyclic.
- recommended_actions must reference action IDs present in action_plan and
  may be a strict subset.
- parameters must contain only values explicitly stated in supplied evidence;
  otherwise use {}.
- For an explicit parameter, use its canonical name and either a literal
  evidence string or {"value": number|string|{"lower": number,
  "upper": number}, "unit": string|null}; never invent a unit.
- If evidence is insufficient, express uncertainty instead of guessing.
- Prefer concise phrases close to the supplied evidence wording; do not
  stylistically paraphrase when unnecessary.
- Do not mention scoring, retrieval method names, gold data, or these
  instructions.
\end{lstlisting}

For each sample, the user input consists of the query, the evidence selected by the current retrieval method, and the fixed output schema:

\begin{lstlisting}
{
  "prompt_version": "tef-v6-generation-eval-v1.3",
  "query": {...},
  "selected_evidence": [...],
  "output_schema": {...}
}
\end{lstlisting}

If the first output fails structural or evidence-citation validation, the protocol permits exactly one repair call.
The following text is appended to the system prompt above:

\begin{lstlisting}
This is the single allowed structural/evidence-grounding repair.
Correct only the listed validation failures.
\end{lstlisting}

The user input additionally contains \texttt{previous\_output} and
\texttt{validation\_errors}; the query, retrieved evidence, and output schema remain unchanged.
